\documentclass[10pt,a4paper]{amsart}
\usepackage[margin=19mm]{geometry}
\usepackage{amsmath,amssymb,mathtools}
\usepackage[scr=rsfs,cal=euler]{mathalfa}
\usepackage{xfrac}
\usepackage[unicode,hidelinks,bookmarks=false]{hyperref}
\newcommand{\ie}{i.e.\ }
\newcommand{\cf}{cf.\ }
\newcommand{\resp}{respectively\ }
\newcommand{\Subsection}{\subsection}
\providecommand{\nobreakdash}{\nobreak\hbox{-}}
\setlength{\emergencystretch}{2em}
\multlinegap0pt
\usepackage{bm}
\usepackage{bbm}
\usepackage{dsfont}
\usepackage{xspace}
\usepackage{cancel}
\usepackage{mathtools}
\usepackage{amsthm}
\makeatletter
\@ifundefined{psmallmatrix}{\newtheorem{psmallmatrix}{Psmallmatrix}}{}
\@ifundefined{theorem}{\newtheorem{theorem}{Theorem}}{}
\@ifundefined{remark}{\newtheorem{remark}[theorem]{Remark}}{}
\makeatother
\title[Introduction to Elliptic Quasi-Modular Forms via Moduli Spac]{Introduction to Elliptic Quasi-Modular Forms via Moduli Spaces}
\author{Walter Andr\'es P\'aez Gaviria}
\date{Source: arXiv:2601.08066v1 (CC BY 4.0)}
\begin{document}
\maketitle

\noindent\emph{Source and licence.} Introduction to Elliptic Quasi-Modular Forms via Moduli Spaces. Version arXiv:2601.08066v1 (CC BY 4.0), distributed under the Creative Commons Attribution 4.0 International licence, as recorded in the arXiv licence metadata for this exact version and shown on the abstract page https://arxiv.org/abs/2601.08066v1. Full rights notice: licenses/arxiv-2601.08066-CC-BY-4.0.txt.\par\smallskip

\noindent\textit{Editorial scope.} Counted unit: the Theorem whose source label is \texttt{bigteo} in the article (source0.tex lines 373--383, source label \texttt{bigteo}), delivered together with the article's own complete original proof of it (source0.tex lines 385--458, one \texttt{proof} environment of 74 lines). No mathematics is written, repaired or re-derived: statement and proof are the article's own text line for line. Required context retained uncounted: remark (lines 296--297). Every definition, display and label of those lines is retained unchanged. Omitted from this package and not redistributed: 1--295, 298--372, 384--384, 459--652. Nothing inside a retained excerpt is omitted or altered: every retained body is delivered exactly as the article's lines are written, so no omission receipt of any kind accompanies this package. Citations: the retained excerpts carry no citation command at all, so no bibliography entry of the article is transcribed and no reference is redistributed. Reference adaptation: none was needed and none was made; every reference inside the delivered unit resolves inside it, so no unresolved reference or citation is produced. Printed slips kept literal: no printed word, symbol, hypothesis or number of the counted statement or of its proof was altered. Layout and class: the article's own class is \texttt{article} with packages including babel, csquotes, amsmath, graphicx, amssymb, amsthm, mathrsfs, authblk, mathtools, biblatex, tikz, bbm, tikz-cd, hyperref; the wrapper is the lane's pinned \texttt{amsart} editorial wrapper at 10pt on A4, which restates the theorem-like environments the retained text needs and adds the article's own macro definitions that the retained text needs (none), with extra wrapper packages (bm, bbm, dsfont, xspace, cancel, mathtools). The delivered numbering is the unit's own. Assets: no retained span contains a figure environment, a graphics-inclusion command, a PDF-page inclusion, a table or a TikZ picture; no image file is copied into this package, the package declares no asset dependency and no image file is redistributed with it. Licence: version arXiv:2601.08066v1 of this article is distributed under the Creative Commons Attribution 4.0 International licence, as recorded in the arXiv licence metadata for this exact version and shown on the abstract page https://arxiv.org/abs/2601.08066v1. A full rights notice is delivered at licenses/arxiv-2601.08066-CC-BY-4.0.txt. Characters: every retained excerpt is pure ASCII, so the delivered engine, XeLaTeX with its default Latin Modern text font, reports no missing character.
\begin{remark} Since $(E_{t_2,t_3},\begin{psmallmatrix}1 & 0\\  t_1 & 1\end{psmallmatrix}\begin{psmallmatrix}\frac{dx}{y} \\ \frac{xdx}{y}\end{psmallmatrix})\cong (E_{t_1,t_2,t_3},\frac{dx}{y},\frac{xdx}{y})$, where  $E_{t_1,t_2,t_3}$ is given in affine coordinates as the zero-set of the polynomial $y^2-(x-t_1)^3+t_2(x-t_1)+t_3$, we see that the family \begin{equation}(E_{t_1,t_2,t_3},\frac{dx}{y},\frac{xdx}{y})), t_1,t_2,t_2\in\mathbb{C}, \Delta\neq 0\end{equation}is an universal family of enhanced elliptic curves. 
\end{remark}\label{remark}

\begin{theorem}\label{bigteo}
For each $\mathfrak{g}\in \mathfrak{G}$, there exists a unique algebraic vector field $\mathsf{R}_{\mathfrak{g}}\in \Theta_{\mathsf{T}}$ such that
\begin{equation}\label{equation}
\nabla_{\mathsf{R}_{\mathfrak{g}}} \begin{psmallmatrix}\alpha\\
\beta\end{psmallmatrix}=\mathfrak{g}^T\begin{psmallmatrix}\alpha\\
\beta\end{psmallmatrix}.
\end{equation}
Here, $\begin{psmallmatrix}\alpha\\
\beta\end{psmallmatrix}=\begin{psmallmatrix}1 & 0\\  t_1 & 1\end{psmallmatrix}\begin{psmallmatrix}\frac{dx}{y} \\ \frac{xdx}{y}\end{psmallmatrix}$ as in Remark \ref{remark}.

\end{theorem}

\begin{proof}
We are going to explicitly compute the vector fields in the statement of this theorem. Let us fix an arbitrary $\mathfrak{g}\in \mathfrak{G}$, and write $\mathsf{R}=\mathsf{R}_{\mathfrak{g}}$ and $S=\begin{psmallmatrix}1 & 0\\  t_1 & 1\end{psmallmatrix}$. We begin by observing that, by the previous lemma, the existence of an unique holomorphic vector field $\mathsf{R}$ on $\mathsf{T}$ satisfying Equation \ref{equation} is already guaranteed. To check that it is algebraic, we use a linear algebra argument. Let 

\begin{equation}
\mathsf{R}=a_1\frac{\partial}{\partial t_1}+ a_2\frac{\partial}{\partial t_2}+a_3\frac{\partial}{\partial t_3}, 
\end{equation} 
where $a_k$ are holomorphic functions on $\mathsf{T}$. Now, we aim to prove that these functions are indeed regular.

First, a calculation using tame polynomials theory (see Chapter 3) yields

\begin{equation}
\mathsf{GM}_{\begin{psmallmatrix}\frac{dx}{y}\\ \frac{xdx}{y}\end{psmallmatrix}}=
\frac{1}{\Delta}\begin{bmatrix}
-\frac{1}{12}d\Delta & \frac{3}{2} \alpha\\
-\frac{1}{8}t_2\alpha & \frac{1}{12}d\Delta 
\end{bmatrix}, \alpha=3t_3dt_2-2t_2dt_3.
\end{equation}

Since $\mathsf{GM}_{\
\begin{psmallmatrix}\alpha\\\beta
\end{psmallmatrix}
}=(dS+S\mathsf{GM}_{\begin{psmallmatrix}\frac{dx}{y}\\ \frac{xdx}{y}\end{psmallmatrix}})S^{-1}$, Equation \ref{equation} is equivalent to 


\begin{equation}\label{despejar}
\begin{bmatrix}
0 & 0\\
a_1 & 0
\end{bmatrix}
+\frac{1}{\Delta}\begin{bmatrix}
1 & 0\\
t_1 & 1
\end{bmatrix}
\begin{bmatrix}
-\frac{1}{12}(54t_3a_3-3t_2^2a_2) & \frac{3}{2}(3t_3a_2-2t_2a_3)\\
-\frac{1}{8}t_2(3t_3a_2-2t_2a_3) & \frac{1}{12}(54t_3a_3-3t_2^2a_2)
\end{bmatrix}=dS+S\mathsf{GM}_{\begin{psmallmatrix}\frac{dx}{y}\\ \frac{xdx}{y}\end{psmallmatrix}}=\mathfrak{g}^TS.
\end{equation}

Entries (1,1) and (1,2) of the previous equality allow us to produce the following linear system
\begin{equation}
\frac{1}{\Delta}\begin{bmatrix}
\frac{1}{4}t_2^2 & -\frac{27}{6}t_3\\
\frac{9}{2}t_3 & -3t_2
\end{bmatrix}\begin{bmatrix}
a_2 \\
a_3
\end{bmatrix}=\begin{bmatrix}
(\mathfrak{g}^TS)_{11}\\
(\mathfrak{g}^TS)_{12}.
\end{bmatrix}
\end{equation}
It can be explicitly inverted to give us 

\begin{equation}
\begin{bmatrix}
a_2 \\
a_3
\end{bmatrix}=\frac{4}{3}\begin{bmatrix}
-3t_2 & \frac{27}{6}t_3\\
-\frac{9}{2}t_3 & \frac{1}{4}t_2^2
\end{bmatrix}\begin{bmatrix}
(\mathfrak{g}^TS)_{11}\\
(\mathfrak{g}^TS)_{12}
\end{bmatrix}.
\end{equation}

Furthermore, entry (2,1) in Equation \ref{despejar}, gives us 

\begin{equation}a_1=(\mathfrak{g}^TS)_{21}-\frac{1}{\Delta}(-\frac{1}{12}t_1(54t_3a_3-3t_2^2a_2)-\frac{1}{8}t_2(3t_3a_2-2t_2a_3)).
\end{equation}

Therefore, $a_1,a_2,a_3$ are global regular functions on $\mathsf{T}$.
\end{proof}

\end{document}
